% ----------------------
\subsection{An introduction to Climbing the Rainbows}
% ----------------------

	Climbing the Rainbows is a series of essays about Latter-day Saint thought. It will present what I believe are some useful ways of thinking about various doctrinal, theological, and cultural issues. More than anything, the essays are a vehicle for me to clarify and develop some views I've entertained off and on over the years. As such, it's really only a collection of opinions, and should be weighed accordingly.

	One impetus for CtR (I didn't choose the name for the fortuitous acronym, but you take what you can get) is the glut of people I have seen leaving the LDS church in the last few years. From my perspective, these people are almost always making a serious mistake, and are doing so with a very questionable rationale. While I have encountered some exceptions, these church-leavers usually carry very simplistic and naive assumptions about the church, God, themselves, revelation, the scriptures, and other topics. I almost never---and here there are even fewer exceptions---see these people interrogating those assumptions. The result is that they make life-altering decisions for bad reasons. The least you can do is make those decisions for good reasons! One of the purposes of these essays is to explain some of the assumptions I have seen, and why I think they are wrong.

	But that's just a secondary goal. As I said before, the primary purpose of these writings is for me to get clear, for my own sake, about what my own beliefs are, and what grounds they have. This will require ranging fairly widely among many different issues, but that shouldn't be surprising, since one of the most important claims of the restored gospel is its universality, both on and beyond the earth.

	I'll try to present the essays in a rough order of logical priority: preliminaries and foundations first, followed by the application of those principles to particular issues. Some of the essays will be about the logical priority itself, and why I think it is the way it is. If we are not taking a principled approach to particular issues, then we increase our chances of making mistakes. But first we have to know what those principles actually are.
	
	With that in mind, on to the essays.